Visible dark-photon phase space and full-exposure-equivalent HPS projections.
Top: the published engineering-run limits in their world-data context; the dashed
rectangle marks the common zoom window in (a) and (b). Gray shading is the union
of the archived visible exclusions, with their original boundaries retained. The world-data contours include BaBar, KLOE/KLOE-2, A1/MAMI,
APEX, NA48/2, HADES, PHENIX, NA64, LHCb, and the older beam-dump or low-mass
constraints from E137, E141, E774, Orsay, KEK, and U70
\cite{BaBarDarkPhoton2014,KLOE2013,KLOE22015,A1MAMI2014,APEX2011,
NA48DarkPhoton2015,HADESDarkPhoton2014,PHENIXDarkPhoton2015,NA64DarkPhoton2018,
NA64X17Visible2021,LHCbDarkPhoton2020,E137BeamDump1988,E141BeamDump1987,
E774BeamDump1991,OrsayBeamDump1989,KEKBeamDump1986,U70DarkPhoton2014}. The thermal
relic target band is shown as a benchmark for sub-GeV dark matter motivation
\cite{EssigEtAl2011}. The HPS 2015 and 2016 engineering-run prompt limits are
highlighted using the legacy display convention that multiplies the published 95\%
$\eps^2$ contours by $1.64/1.96$ \cite{HPS2015DarkPhoton,HPS2016PromptLong}; the
NA48/2 and BaBar inputs are treated in the same way. This approximate Gaussian-ratio
rescaling is retained only for historical context and is not a recalculation of a
90\% CL upper limit with the \CLs{} construction. The lepton central-value lines use $\Delta a_\mu=38\times10^{-11}$
(WP25) and $\Delta a_e\simeq0.34\times10^{-12}$ (Rb20 with Fan23)
\cite{MuonWP2025,ElectronFan2023,AlphaMorel2020}. These are reference loci,
not confidence bands; WP25 is compatible with zero.
(a) Full-equivalent 2015+2016+2021; (b) full-equivalent 2015+2016+2019+2021.
The latter adds the measured 2019 nominal 1\% spectrum as a separate likelihood
input with a 2021-response proxy. Both current-sample combined limits are then scaled by the square root
of the current-to-full summed count-density ratio, taking 2021 from 10\% to
100\% and, in (b), 2019 from 1\% to 100\%. These are conditional
observed-equivalent projections, not expected sensitivities or full-data
exclusions; Appendix~\ref{app:figure2-projections} gives the construction.
